\editionPage[\pdfbookmark[0]{Fiches outils}{bm-outils-debut-parent}]{outils-debut}{Fiche outil}{Calculs et algèbre}{À consulter au fil des chapitres.}
\begin{blocEdition}{outils-debut}

\editionEtape{Relatifs et fractions}Prévoir le signe d’un produit ou d’un quotient avant de calculer. Pour additionner des fractions, utiliser un même dénominateur. Pour multiplier, multiplier les numérateurs entre eux et les dénominateurs entre eux. Pour diviser par un nombre non nul, multiplier par son inverse.
\editionEtape{Puissances et racines}$a^n$ est le produit de $n$ facteurs égaux à $a$. $a^m\times a^n=a^{m+n}$ et $a^n\times b^n=(ab)^n$. Pour $b$ positif, $\sqrt b$ est le nombre positif dont le carré vaut $b$ ; on l’encadre avec deux carrés parfaits.
\editionEtape{Distributivité}$k(a+b)=ka+kb$ et $k(a-b)=ka-kb$. Développer va du produit vers la somme ; factoriser fait le chemin inverse. $-(a-b)=-a+b$. Réduire seulement des termes de même nature.
\editionEtape{Équation}Développer et réduire, effectuer la même opération sur les deux membres, isoler l’inconnue, puis vérifier dans l’équation initiale et dans le contexte.
\editionRetenir{Une égalité exacte s’écrit avec $=$, une valeur arrondie avec $\approx$. Conserver le résultat exact jusqu’au dernier calcul.}
\end{blocEdition}
\editionPage{outils-2}{Fiche outil}{Proportionnalité et grandeurs}{À consulter au fil des chapitres.}
\begin{blocEdition}{outils-2}

\editionEtape{Proportionnalité}Coefficient, linéarité, retour à l’unité ou produits en croix : pour $b\ne0$ et $d\ne0$, $\frac ab=\frac cd$ équivaut à $ad=bc$. Une proportionnalité se représente par des points alignés avec l’origine.
\editionEtape{Ratios et partages}Partager selon le ratio $a:b:c$ : compter $a+b+c$ parts, calculer une part, puis multiplier. L’ordre des nombres d’un ratio compte.
\editionEtape{Pourcentages}Hausse de $t\,\%$ : multiplier par $1+\frac t{100}$. Baisse de $t\,\%$ : multiplier par $1-\frac t{100}$. Évolutions successives : multiplier les coefficients.
\editionEtape{Grandeurs quotients}$v=\frac dt$, $q=\frac Vt$, $\rho=\frac mV$. $1$ h $=3\,600$ s ; de km/h à m/s, diviser par 3,6. $1$ m$^3=1\,000$ L et $1$ L $=1$ dm$^3$.
\editionEtape{Aires et volumes}Triangle : $\frac{b\times h}2$. Disque : $\pi r^2$. Prisme et cylindre : $\mathcal B\times h$. Pyramide et cône : $\frac13\mathcal B\times h$.
\editionRetenir{Avant de calculer avec deux mesures, choisir des unités compatibles et écrire l’unité du résultat.}
\end{blocEdition}
\editionPage{outils-3}{Fiche outil}{Choisir un outil géométrique}{À consulter au fil des chapitres.}
\begin{blocEdition}{outils-3}

\editionEtape{Droites remarquables}M est sur la médiatrice de $[AB]$ si et seulement si $MA=MB$. Hauteur : perpendiculaire issue d’un sommet. Médiane : sommet et milieu du côté opposé ; elle partage l’aire en deux.
\editionEtape{Translations et parallélogrammes}Une translation conserve longueurs, angles et aires ; si A a pour image A’ et B pour image B’, ABB’A’ est un parallélogramme. Diagonales de même milieu : parallélogramme ; de plus de même longueur : rectangle ; perpendiculaires : losange.
\editionEtape{Des milieux dans un triangle}Deux milieux : parallélisme et moitié du troisième côté. Un milieu et une parallèle : un nouveau milieu.
\editionEtape{Triangle rectangle}Pythagore pour une longueur ; réciproque pour prouver un angle droit ; contraposée pour prouver qu’il n’y en a pas. Le centre du cercle circonscrit est le milieu de l’hypoténuse ; un triangle inscrit dans un demi-cercle est rectangle.
\editionRetenir{Une preuve cite ses données et la propriété utilisée. Un codage de l’énoncé est une donnée ; une mesure sur le dessin n’est qu’un indice.}
\end{blocEdition}
\editionPage{outils-4}{Fiche outil}{Données, hasard et programmation}{À consulter au fil des chapitres.}
\begin{blocEdition}{outils-4}

\editionEtape{Statistiques}Moyenne pondérée : somme des produits valeur $\times$ effectif, divisée par l’effectif total. Médiane : ranger, puis partager la série en deux parties de même effectif. Étendue : maximum moins minimum. Tableur : MOYENNE, MEDIANE, MAX, MIN, NB.
\editionEtape{Probabilités}Écrire l’univers, puis l’événement comme ensemble d’issues. $P(\overline A)=1-P(A)$. $A\cap B$ : issues communes ; $A\cup B$ : issues d’au moins un des deux. Deux épreuves : arbre ou tableau à double entrée.
\editionEtape{Fréquences}Une fréquence observée fluctue d’une série à l’autre ; elle donne une approximation de la probabilité, sans la remplacer.
\editionEtape{Scripts}Repérer le déclencheur, initialiser la variable, suivre les blocs dans l’ordre, compter les tours, tester les cas limites d’une condition. Dans un tableur, une référence relative change quand on recopie une formule.
\editionRetenir{Après correction, refaire la question sans aide et expliquer l’idée qui l’a résolue : le but est de retrouver la méthode.}
\end{blocEdition}
